\section{Related Work}
\label{sec:literature}

This section situates \dreamprice{} within seven strands of literature: traditional demand forecasting and price optimization, learned world models for sequential decision making, structured state-space models, causal inference in demand estimation, offline reinforcement learning, pricing optimization and retail analytics, and the Dominick's scanner dataset as the canonical empirical setting.

\subsection{Traditional Demand Forecasting and Price Optimization}

The standard approach to retail pricing proceeds in two stages: estimate demand as a function of price and other covariates, then optimize prices subject to business constraints. On the forecasting side, classical time-series methods, including autoregressive integrated moving average (ARIMA) models, exponential smoothing, and seasonal decomposition, have long served as workhorses for retail demand prediction \citep{fildes2022retail}. These methods model temporal patterns in aggregate sales but do not capture cross-product substitution or the causal effect of price changes on demand. Machine learning methods, including gradient boosting machines \citep{friedman2001greedy}, random forests, and neural networks, have improved point forecast accuracy by incorporating high-dimensional features such as promotions, holidays, competitor prices, and weather. \citet{fildes2022retail} provide a comprehensive survey of retail forecasting methods, documenting that ML approaches now dominate industry practice for point forecasting while statistical methods remain competitive for uncertainty quantification.

On the optimization side, the dominant paradigm is to estimate price elasticities, typically via log-linear demand models or the Berry-Levinsohn-Pakes (BLP) framework \citep{berry1995automobile}, and then solve a constrained optimization problem. Common formulations include linear or quadratic programming to maximize gross margin subject to price-range constraints, competitive parity rules, and promotional calendars \citep{talluri2004revenue}. Ramsey pricing, which sets markups inversely proportional to demand elasticity, provides a welfare-theoretic benchmark. These two-stage approaches are transparent and interpretable but fundamentally limited: the demand model is static (fit once, applied repeatedly), the optimization horizon is typically a single period, and the system cannot capture multi-step dynamics such as stockpiling, reference-price effects, or competitive responses that unfold over weeks. \dreamprice{} replaces both stages with an end-to-end learned dynamics model that supports multi-step imagination and policy optimization over arbitrary horizons.

\subsection{World Models for Sequential Decision Making}

The idea of learning a world model, a compact representation of environment dynamics that supports prediction and planning, dates to \citet{ha2018worldmodels}, who demonstrated that a variational autoencoder could learn a latent space in which an agent could imagine future trajectories. That foundational work showed that an agent could learn to navigate a maze by training purely in imagined latent space, without ever observing the true environment during policy learning. The subsequent decade has seen a steady progression of architectures that refine this idea.

The Dreamer family \citep{hafner2020dreamerv1, hafner2021dreamerv2, hafner2023dreamerv3} established the recurrent state-space model (RSSM) as the dominant paradigm for learning world models. DreamerV1 \citep{hafner2020dreamerv1} introduced latent imagination: the agent learns a stochastic latent representation that evolves over time, with a decoder predicting observations and rewards from the latent state. DreamerV2 \citep{hafner2021dreamerv2} extended this to discrete action spaces by mastering Atari games through imagination-based policy optimization, demonstrating that world models could scale to high-dimensional pixel observations. DreamerV3 \citep{hafner2023dreamerv3} unified the design across discrete and continuous domains, introducing symlog normalization, twohot distributions for rewards and values, and KL balancing to prevent posterior collapse. The RSSM architecture, a deterministic recurrent core (traditionally a GRU) augmented with stochastic latent variables, has become the standard for model-based reinforcement learning in physical environments.

Parallel developments have explored alternative architectures. MuZero \citep{schrittwieser2020muzero} demonstrated that planning with a learned model suffices for mastering board games and Atari, using a learned dynamics model that predicts value and policy targets rather than raw observations. The model is trained end-to-end with Monte Carlo tree search, and the resulting system achieves superhuman performance in Go, chess, and shogi. IRIS \citep{micheli2023iris} and TD-MPC2 \citep{hansen2024tdmpc2} have extended world models to continuous control and high-dimensional observations. IRIS replaces the RSSM with a transformer-based architecture, using discrete tokens for the latent representation and achieving strong performance on DMC and Atari. TD-MPC2 combines model-based planning with temporal difference learning for continuous control, scaling to complex robotics tasks.

A critical observation unifies these works: they all operate on physical environments. The Dreamer family, MuZero, IRIS, and TD-MPC2 are evaluated on games (Atari, board games), physical simulators (MuJoCo, Isaac Gym), or robotic control. The dynamics they learn are governed by physics or game rules: the transition from state to state is exogenous to the agent's policy in the sense that the laws of motion do not change with the agent's behavior. In economic environments, by contrast, the transition dynamics are endogenous. Prices and demand are simultaneously determined; retail prices respond to expected demand, costs, and competitive pressure, while demand responds to those prices. No learned world model has been applied to such economic domains. \dreamprice{} fills this gap by learning the first retail dynamics model from scanner data.

\subsection{Structured State Space Models}

The recurrent backbone of standard RSSMs is typically a gated recurrent unit (GRU) or long short-term memory (LSTM) network. These architectures process sequences sequentially, with each time step depending on the previous hidden state. For long sequences, this sequential dependency can be a bottleneck; the GRU cannot exploit parallel computation across the time dimension during training.

Structured state space models (SSMs) offer an alternative that combines long-range dependency modeling with computational efficiency. \citet{gu2022s4} introduced S4 (Structured State Spaces), which parameterizes sequences using linear state-space equations with learned transition matrices. S4 achieves strong performance on long-range dependencies, including image classification and audio generation, by using the HiPPO initialization for the state matrix and efficient parallel scan algorithms. The key insight is that linear recurrence admits a closed-form solution that can be computed in parallel via associative scan operations.

Mamba \citep{gu2024mamba} extended S4 with selective gating: the state transition parameters become input-dependent rather than fixed. This selectivity allows the model to focus on relevant information at each step, analogous to the gating in attention mechanisms. Mamba achieves linear-time complexity in sequence length and scales to long contexts without the quadratic cost of attention. Mamba-2 \citep{dao2024mamba2} further unifies these ideas through structured state-space duality (SSD), which establishes a formal connection between linear attention and state-space models. The SSD framework shows that both can be viewed as instances of a more general class of models, enabling algorithms that share the best properties of each: parallel training with $O(n)$ complexity in sequence length, and recurrent inference with $O(1)$ cost per step.

This duality has direct implications for world models. During training, the model processes sequences of observations and actions in parallel: the entire batch of sequences can be scanned in one pass, as the posterior over the stochastic latent depends only on the current observation in the DRAMA-style decoupled design \citep{rajbhandari2024drama}. During imagination, the model switches to recurrent mode: the Mamba-2 block's \texttt{step()} function computes the next hidden state from the previous state and the current input in constant time per step. This matches the training-then-imagination workflow of DreamerV3, but with a backbone that scales more efficiently to long sequences and avoids the sequential bottleneck of GRU recurrence. \dreamprice{} adopts this architecture, replacing the GRU with Mamba-2 for the recurrent core.

\subsection{Causal Inference in Demand Estimation}

Estimating demand as a function of price is a central problem in empirical industrial organization and marketing. The naive approach of regressing quantity on price fails because prices are endogenous. Retailers set prices in response to expected demand, costs, and competitive conditions; unobserved factors that correlate with both price and demand (e.g., product quality, seasonality, local shocks) confound the regression. The resulting elasticity estimates are biased, and a world model built on such estimates would produce misleading counterfactuals.

The econometric literature has developed several identification strategies. The Berry-Levinsohn-Pakes (BLP) framework \citep{berry1995automobile} uses a structural model of discrete choice with random coefficients; identification relies on moment conditions that match predicted choice shares to observed shares. \citet{nevo2001measuring} applied this approach to the ready-to-eat cereal industry, estimating elasticities and measuring market power. \citet{hausman1978specification} introduced the use of instrumental variables for specification testing and parameter identification. When prices are correlated across stores or time for reasons exogenous to store-specific demand shocks, cross-sectional or temporal variation can serve as instruments. The Hausman IV approach, which uses the average log price in other stores as an instrument for a given store's price, exploits the panel structure of scanner data: prices in other stores are relevant (drivers of cost and competition are correlated) but plausibly exogenous to store-specific demand shocks.

\citet{hoch1995determinants} applied this logic to the Dominick's dataset, estimating store-level price elasticities using cross-store instruments. Their work demonstrated that elasticity varies substantially across stores and product categories, with important implications for targeted pricing. \citet{park2010copula} proposed a nonparametric copula-based approach to demand estimation with endogeneity, using the joint distribution of price and control variables to construct a control function. The copula correction does not require specifying instruments ex ante but relies on distributional assumptions.

The double machine learning (DML) framework of \citet{chernozhukov2018dml} provides a flexible approach that can incorporate any of these identification strategies. DML partitions the estimation into orthogonal nuisance parameters and parameters of interest, using cross-fitting to avoid overfitting bias. The framework accommodates high-dimensional controls and nonparametric first-stage estimators. \citet{bach2022doubleml} provide a Python implementation that supports DML-PLIV (partial linear instrumental variables) for demand estimation with endogenous prices.

The key contribution of \dreamprice{} in this literature is the integration of econometric identification into a neural world model. Rather than estimating demand entirely from data, which would conflate the causal effect of price with confounders, \dreamprice{} pre-estimates causal price elasticities via DML-PLIV and freezes them into the demand decoder. The decoder's residual MLP learns the remainder, including seasonality, demographics, promotion effects, and substitution patterns, while the price-demand relationship remains causally identified. Freezing econometric estimates into neural network decoders is novel; prior work in demand estimation either uses parametric structural models or applies machine learning without causal constraints.

\subsection{Offline Reinforcement Learning}

Reinforcement learning methods can be classified along two orthogonal axes: the data regime (online versus offline) and the use of a dynamics model (model-free versus model-based). Online methods interact with the environment during training, collecting new experience and refining their policy through exploration. Model-free online methods, such as DQN \citep{mnih2015dqn}, PPO \citep{schulman2017ppo}, and SAC \citep{haarnoja2018sac}, learn value functions or policies directly from transitions without building an explicit dynamics model. Model-based online methods, including the Dreamer family \citep{hafner2023dreamerv3} and MuZero \citep{schrittwieser2020muzero}, learn a dynamics model and optimize the policy in imagination, reducing the number of real-world interactions required.

In many settings, retail pricing among them, online exploration is infeasible or prohibited: experimenting with prices on real customers carries financial risk and ethical constraints. The agent must learn from a fixed dataset of historical trajectories collected under some (possibly unknown) behavioral policy. This offline learning setting introduces the distribution shift problem: the policy being optimized may visit states and actions that differ from the data distribution, and the learned value function or dynamics model may extrapolate poorly outside the support of the data \citep{levine2020offline}.

Model-free offline methods address distribution shift through value-function conservatism or policy constraints. Batch-Constrained Q-learning (BCQ) \citep{fujimoto2019bcq} restricts the policy to actions observed in the dataset by using a generative model of the behavioral policy. Conservative Q-Learning (CQL) \citep{kumar2020cql} penalizes the Q-function for assigning high values to out-of-distribution actions, effectively learning a lower bound on the true Q-function. Implicit Q-Learning (IQL) \citep{kostrikov2022iql} avoids querying out-of-distribution actions entirely by using expectile regression to approximate the optimal value function from the dataset. Decision Transformer \citep{chen2021decision} reframes offline RL as a sequence prediction problem, conditioning a transformer on desired returns to generate actions. These methods learn policies but cannot perform counterfactual reasoning or multi-step imagination.

Model-based offline methods combine a learned dynamics model with pessimism to handle distribution shift. MOReL \citep{kidambi2020morel} constructs a pessimistic uncertainty set around the learned dynamics model and restricts planning to states that are provably within the uncertainty set. MOPO \citep{yu2020mopo} takes a softer approach: it modifies the reward signal during imagination rather than constraining the policy or dynamics. The key idea is that the learned dynamics model is uncertain in regions of state-action space that are poorly covered by the data. MOPO uses an ensemble of dynamics models to estimate this uncertainty; the reward is penalized by the uncertainty, so the agent is discouraged from visiting states where the model is unreliable. The penalty is $r_{\text{pessimistic}} = r_{\text{mean}} - \lambda_{\text{LCB}} \cdot r_{\text{std}}$, where $r_{\text{mean}}$ and $r_{\text{std}}$ come from the ensemble. COMBO \citep{yu2021combo} combines model-based and conservative approaches, using offline data to augment the model rollout and applying conservative regularization to the value function.

\dreamprice{} occupies the model-based offline quadrant: it learns a world model from historical scanner data and trains a policy entirely in imagination, using MOPO-style pessimism to guard against distributional shift. The MOPO approach is particularly well-suited to the DreamerV3 training loop. DreamerV3 already trains an actor in imagination; the only modification is to change the reward signal used during actor training. The critic and world model remain unchanged. A 5-head reward ensemble provides $r_{\text{mean}}$ and $r_{\text{std}}$; the actor is trained on the pessimistic reward while the critic continues to predict the true return. This minimal modification preserves the three-phase structure and avoids the complexity of full conservative value estimation.

\subsection{Pricing Optimization and Retail Analytics}

The application of reinforcement learning to pricing has a substantial literature. \citet{chen2022dynamic} study dynamic pricing with reinforcement learning, formulating the problem as a Markov decision process and analyzing the trade-off between exploration and exploitation. \citet{ban2021personalized} extend this to personalized pricing with heterogeneous elasticity, using machine learning to estimate demand as a function of high-dimensional features. The AI Economist \citep{zheng2020aieconomist} applies multi-agent reinforcement learning to tax policy design, training agents in a rule-based economic environment to learn tax schedules that balance equality and productivity.

These works share a common limitation: they rely on structural simulators or model-free methods. Chen and Simchi-Levi assume a known or estimated demand model; Ban and Keskin use model-free RL with feature-based demand estimation. The AI Economist uses a rule-based economic environment for training. None of these approaches learns a dynamics model from data. The gap is twofold: (1) no learned world model exists for pricing environments, and (2) no imagination-based planning has been applied to retail. A model-free agent can learn a policy from historical data, but it cannot answer counterfactual queries such as ``what if the price of SKU 3 were cut by 5\% next quarter?'' without executing the change. A learned world model enables such queries by rolling out latent trajectories under hypothetical actions.

\dreamprice{} is the first system to learn a dynamics model for retail pricing from scanner data. The model captures demand responses to prices, promotions, and seasonality; it supports imagination-based policy optimization; and it integrates causal identification to prevent confounded elasticity estimates. The architecture is designed for the structured economic state space (prices, quantities, demographics, and derived features in 100--300 dimensions), which is fundamentally more tractable than the high-dimensional pixel observations that DreamerV3 was designed for.

\subsection{Scanner Data and the Dominick's Dataset}

The Dominick's Finer Foods dataset \citep{dominicks1999} has become the canonical dataset for empirical retail and marketing research. Maintained by the James M. Kilts Center for Marketing at the University of Chicago Booth School of Business, the dataset spans 400 weeks (September 1989 to May 1997), 93 stores, and approximately 18,000 unique UPCs across 29 product categories. The movement files record store-SKU-week tuples with unit sales, prices, profit margins, and promotion indicators. The data structure enables rich panel analysis: the same products are observed across stores and over time, with variation in prices driven by local competition, cost shocks, and promotional calendars.

\citet{hoch1995determinants} used the Dominick's data to estimate store-level price elasticities, demonstrating that elasticity varies substantially across stores and that store demographics (income, education, competition) explain much of this variation. \citet{montgomery1997estimating} estimated heterogeneous price thresholds, the price levels at which consumers switch between brands, using the same dataset. The panel structure (93 stores) supports Hausman-style instruments: the average log price in other stores for a given UPC-week serves as an instrument for a store's own price, as it is correlated with cost and competition but plausibly exogenous to store-specific demand shocks. The store demographic file, with 300+ variables from the 1990 Census, enables control for heterogeneity in consumer preferences and competitive intensity.

The Dominick's dataset is ideal for training a learned world model. It directly contains state-action-outcome tuples: the action is the price and promotion decision, the state includes competitor context, seasonality, and store demographics, and the outcome is observed demand. The PROFIT field enables wholesale cost derivation. The temporal split (train weeks 1--280, validation 281--340, test 341--400) is strictly chronological, avoiding look-ahead bias. Categories such as canned soup, beer, and soft drinks offer moderate SKU counts with clear substitution patterns and promotional activity. \dreamprice{} is trained on this dataset, with preprocessing that constructs Hausman IVs, computes derived features, and applies symlog normalization for the world model encoder.
